\documentclass[12pt]{article}
\usepackage[margin=1in]{geometry}
\usepackage{amsmath,amsfonts,amssymb}
\usepackage{graphicx}

\title{Derivation of Control System for a 3-DOF Quadcopter State Space Model}
\date{Spring 2024}

\begin{document}

\maketitle
\section{Introduction}
This project document outlines the theoretical framework and project requirements for designing a control system for a quadcopter. It includes the basics of quadcopter dynamics, state space representation, suggested control methods with resources, and a MATLAB code template to get started!

\section{Quadcopter Dynamics \& Simplifications}
The quadcopter's movement is controlled through forces and torques generated by its four rotors. This section provides an overview of the fundamental dynamics of a quadcopter.

\subsection{Roll, Pitch, and Yaw}
\begin{itemize}
    \item \textbf{Roll}: Rotation about the quadcopter's longitudinal axis. Controlled by creating a differential thrust between the left and right pairs of rotors.
    \item \textbf{Pitch}: Rotation about the lateral axis. Controlled by the differential thrust between the front and rear rotors.
    \item \textbf{Yaw}: Rotation about the vertical axis. Controlled by varying the relative speeds of rotors spinning in opposite directions.
\end{itemize}

\subsection{General Equation for Thrust and Torque}
Each rotor contributes to the overall thrust and torque of the quadcopter:
\begin{itemize}
    \item Thrust generated by rotor $i$: $T_i = k w_i^2$
    \item Torque generated by rotor $i$: $\tau_i = b w_i^2$
\end{itemize}
where $k$ is the thrust constant, $b$ is the drag coefficient, and $w_i$ is the speed of rotor $i$.

\subsection{Net Forces and Torques}
The net thrust and torques in the roll ($\phi$), pitch ($\theta$), and yaw ($\psi$) directions are given by:
\begin{align*}
    T &= \sum T_i = k(w_1^2 + w_2^2 + w_3^2 + w_4^2) \\
    \tau_{\phi} &= (T_2 + T_4 - T_1 - T_3) \cdot L \\
    \tau_{\theta} &= (T_3 + T_4 - T_1 - T_2) \cdot L \\
    \tau_{\psi} &= \tau_1 - \tau_2 + \tau_3 - \tau_4
\end{align*}
where $L$ is the distance from each rotor to the center of the quadcopter.

\subsection{Euler’s Equations for Rotational Motion}
The rotational dynamics of the quadcopter are governed by Euler's equations:
\begin{align*}
    I_x \ddot{\phi} &= \tau_{\phi} - (I_y - I_z) \dot{\theta} \dot{\psi} \\
    I_y \ddot{\theta} &= \tau_{\theta} - (I_z - I_x) \dot{\phi} \dot{\psi} \\
    I_z \ddot{\psi} &= \tau_{\psi} - (I_x - I_y) \dot{\phi} \dot{\theta}
\end{align*}
where $I_x$, $I_y$, and $I_z$ are the moments of inertia about the roll, pitch, and yaw axes, respectively.

\subsection{Linearization for State Space Model}
For small angle approximations, the sin and cos of the angles are approximated as follows:
\begin{align*}
    \sin(\theta) &\approx \theta \\
    \cos(\theta) &\approx 1
\end{align*}
This approximation simplifies Euler's equations for small angles.

\section{State Space Representation}
The state space model of the quadcopter is represented as:
\begin{align*}
    \dot{X} &= AX + BU \\
    Y &= CX + DU
\end{align*}
where $X$ is the state vector, $U$ is the input vector, and $Y$ is the output vector. The matrices $A$, $B$, $C$, and $D$ are defined based on the system dynamics.

\subsection{State Vector}
The state vector $X$ consists of the angular positions and velocities:
\begin{equation*}
    X = [\phi, \dot{\phi}, \theta, \dot{\theta}, \psi, \dot{\psi}]^T
\end{equation*}

\subsection{Input Vector}
The input vector $U$ consists of the speeds of the four motors:
\begin{equation*}
    U = [w_1, w_2, w_3, w_4]^T
\end{equation*}

\subsection{System Matrix (A)}
The system matrix $A$ in the state space representation is:
\begin{equation*}
    A = \begin{bmatrix}
     
    \end{bmatrix}
\end{equation*}

\subsection{Input Matrix (B)}
The input matrix $B$ relates the motor speeds to the system dynamics:
\begin{equation*}
    B = \begin{bmatrix}
    0 & 0 & 0 & 0 \\
    \frac{-L \cdot k}{I_x} & \frac{L \cdot k}{I_x} & \frac{L \cdot k}{I_x} & \frac{-L \cdot k}{I_x} \\
    0 & 0 & 0 & 0 \\
    \frac{-L \cdot k}{I_y} & \frac{-L \cdot k}{I_y} & \frac{L \cdot k}{I_y} & \frac{L \cdot k}{I_y} \\
    0 & 0 & 0 & 0 \\
    \frac{b}{I_z} & \frac{-b}{I_z} & \frac{b}{I_z} & \frac{-b}{I_z}
    \end{bmatrix}
\end{equation*}

\subsection{Output Matrix (C)}
The output matrix \( C \) in the state space representation defines how the state variables are mapped to the output variables. In many control problems, the output variables are a subset of the state variables or a linear combination thereof.\newline

For our quadcopter system, if we are interested in observing all the state variables (angular positions and velocities), the output matrix \( C \) would be an identity matrix. However, if only angular positions or velocities are of interest, \( C \) would be constructed accordingly.\newline

Assuming we want to observe all the state variables, the output matrix \( C \) is a 6x6 identity matrix:
\begin{equation*}
    C = \begin{bmatrix}
    1 & 0 & 0 & 0 & 0 & 0 \\
    0 & 1 & 0 & 0 & 0 & 0 \\
    0 & 0 & 1 & 0 & 0 & 0 \\
    0 & 0 & 0 & 1 & 0 & 0 \\
    0 & 0 & 0 & 0 & 1 & 0 \\
    0 & 0 & 0 & 0 & 0 & 1
    \end{bmatrix}
\end{equation*}

\subsection{Feedforward Matrix (D)}
The feedforward matrix \( D \) in the state space representation directly relates the input to the output, bypassing the state dynamics. For many physical systems, especially those like our quadcopter model, there is no direct feedforward path from the input to the output. Therefore, in such cases, the feedforward matrix \( D \) is typically a zero matrix.\newline

For our quadcopter system, assuming no direct feedforward effect, the feedforward matrix \( D \) is a zero matrix of appropriate dimensions. If our system has 6 outputs and 4 inputs, \( D \) will be a 6x4 zero matrix:
\begin{equation*}
    D = \begin{bmatrix}
    0 & 0 & 0 & 0 \\
    0 & 0 & 0 & 0 \\
    0 & 0 & 0 & 0 \\
    0 & 0 & 0 & 0 \\
    0 & 0 & 0 & 0 \\
    0 & 0 & 0 & 0
    \end{bmatrix}
\end{equation*}


% Include the State Space Representation Section Here
% Include the System and Input Matrices Section Here

\section{Suggested Control Systems and Resources}
Students are encouraged to explore various control systems. Here are some recommendations:

\begin{itemize}
    \item \textbf{PID Control}: Chapters 4-6 of ``Feedback Control of Dynamic Systems" by Franklin et al.
    \item \textbf{Linear Quadratic Regulator (LQR)}: ``Optimal Control Theory: An Introduction" by Donald E. Kirk.
    \item \textbf{Model Predictive Control (MPC)}: ``Model Predictive Control System Design and Implementation Using MATLAB" by Wang.
    \item \textbf{Other Methods}: Feel free to do your own research into different control methods!
\end{itemize}



\section{Project Requirements}
Students are required to design a control system that meets the following specifications:

\begin{enumerate}
    \item Set the rotational motions (yaw, pitch, roll) to 5 degrees initially.
    \item At t=2s, change yaw to 30 degrees, pitch to 20 degrees, and roll to 0.
    \item At t=10s, set all variables to 0.
    \item Ensure no steady-state error and a rise time of 0.5 seconds per degree.
    \item Plot yaw, pitch, and roll on one graph, and motor inputs(control variables) on another.
    \item Analyze the system's robustness and responsiveness to external disturbances.
\end{enumerate}

Students must \textbf{also} submit a comprehensive report that elucidates the functionality of their code, describes the control system utilized, and explains the tuning process if relevant. The report should also incorporate graphs generated by the MATLAB code, accompanied by detailed explanations of the significance and interpretation of each plot.

\section{Evaluation Criteria}
The project will be evaluated based on the following criteria:
\begin{itemize}
    \item [-] [30\%] Accuracy in achieving the specified rotational motions.
    \item [-] [20\%] Adherence to performance metrics (steady-state error, rise time).
    \item [-] [20\%] Quality of plots and data presentation.
    \item [-] [10\%] Analysis of system robustness and responsiveness.
    \item [-] [20\%] Overall quality of the report and code.
\end{itemize}

\section{MATLAB Code Template}
A MATLAB code template is provided for initial setup. Students will complete the implementation based on their chosen control method.

\begin{verbatim}
% Define the quadcopter parameters
% Define the state space matrices A, B, C, D
% Simulate the system response to control inputs
% Plot the response of yaw, pitch, and roll
% Plot the motor input speeds
% Add your control system implementation here
\end{verbatim}

\end{document}
